% -----------------------------------------------------------
\section{1877}
% -----------------------------------------------------------

	% ----------------------
	\subsection{Philo Dibble}
	% ----------------------
		\label{EMS-1877-PD}

		% ----------------------
		\subsubsection{Introduction to 1877}
		% ----------------------
		
			sdf
			
		% -----
		\subsubsection{Sermon, 7 January 1877}
		% -----
			\label{EMS-1877-PD-7-JAN-1877}
			% \label{EMS-18XX-BY-DAY-3LETTERMONTH-YEAR}
			
			\begin{description}
					\item[Print Sources]
				\end{description}

					\begin{tabular}{p{0.5in}p{1in}p{0.5in}p{1in}}
					---	& ---		& --- 		& --- \\
					\end{tabular}
				
				\begin{description}
					\item[Online Access] \href{https://catalog.churchofjesuschrist.org/assets/d082597d-fef5-4c5e-8741-05eafb08df4a/0/0}{Here}
					% \item[Analysis] \hyperref[XX]{Here}
				\end{description}
				
				\begin{description}
					\item[Background]
				\end{description}
				
					\phantom{.}
				
				\begin{description}
					\item[Original Source]
				\end{description}
				
					Payson Ward general minutes; this is the complete transcript of what Dibble said that day.
					% brief description of original source material, with reference to JSP or somewhere else for more info
					% Background material: it might be helpful to have a note that says whether a given document was presented as a speech, was written in a journal, etc.
					% Also a comment here about how reliable the source might be, or what shape it is in, if there are large omissions or parts that are hard to understand, and so on.
					% Also a comment about when I transcribed it, whether I've checked it more than once, and so on
				
				\begin{description}
					\item[Text]
				\end{description}
				
					\uline{Elder} \uline{Philo} \uline{Dibble} of Springville next addressed the Congregation: Said he was a stranger to most of this congregation, but was not a stranger to the gospel, He had been in the church over 47 years. He was among the first 50 that embraced the gospel in this dispensation. He was present when Jos. Smith and Sidney Rigdon when they had that glorious vision of the creation \&c. It was shown unto them that a hundred and forty four thousand should stand on the earth in the last days as Saviors of men. We all agreed in the heavens to come onto this earth and pass through this probation having forgotten all that we knew before. We are here for trial and they who hold out faithful will receive the crown. He was familiar with the Prophet Joseph from boyhood and knew him to be a man of God. Never loved any one as he loved Joseph. The Lord told the prophet when this church was organized with six members, that he was laying the foundation of a great work. It has already grown to greater dimensions than man could have dreamed of. He was present when the heavens were opened and men saw within the veil and saw the man of sin revealed and many other great and glorious things.
					
					Related many other things that transpired in the Kirtland temple at its dedication. Testified that all that had been predicted by Joseph Smith would be literally fulfilled. The speaker prophesied that more would transpire \uline{this} year than had in the three last, and the same in proportion every succeeding year until the winding up scene.
					
					Referred to the persecutions that the Saints have passed through and the refusal of our enemies to redress our wrongs, how even rulers and men in high places turned a deaf ear to our cries. He enlisted in the cause of Christ as a soldier during the war. He hoped never to turn his back on his brethren: that he would be a valiant soldier to the end, No man who has the spirit of this work ever murmurs. As great and as many miracles had been worked in favor of Israel in our day as ever were in any former day or dispensation. Related some of these miracles and showed how the Lord had favored this people from the begining.
					The signs of the times that are now being made manifest should cause the Saints to lift up their hearts and rejoice. The return of the lost tribes is near at hand---the time to favor Zion. All the signs of the last days foretold by the Savior are before us. This is a great work and the laborers in the field are few. Let out hearts be glad that we live here in peace while the judgments of God are stalking through the earth. Related a dream he had three years ago. From this dream did not believe that either Tilden or Hayes would dare to be inaugurated. --- Said he had been into the spirit world twice and he knew something about it. He saw no children in that world. Believed that spirits of little children do not go into the spirit world, but that being pure and spotless that they go back immediately into the presence of the Father and Son. He saw there that if a man here offend a brother without a cause and do not make it right here in the flesh that sin follows after the perpetrator and he has it to settle there; but that if we settle our difficulties here and make things right, it remains right and does not pursue us beyond this life.
					
					Said he was traveling again with his paintings illustration of the early history of the Church and would exhibit them here some evening for the benefit of the people. He had been called to this workd; it is his mission to keep these things before the saints, that they may not be forgotten. 

	% ----------------------
	\subsection{Brigham Young}
	% ----------------------
		\label{EMS-1877-BY}

		% ----------------------
		\subsubsection{Introduction to 1877}
		% ----------------------
		
			sdf
			
		% -----
		\subsubsection{Instructions, 7 February 1877}
		% -----
			\label{EMS-1877-BY-7-FEB-1877}
			% \label{EMS-18XX-BY-DAY-3LETTERMONTH-YEAR}
			
			\begin{description}
					\item[Print Sources]
				\end{description}

					\begin{tabular}{p{0.5in}p{1in}p{0.5in}p{1in}}
					CDBY	& ????		& JD 		& ???? \\
					\end{tabular}
				
				\begin{description}
					\item[Online Access] \href{https://archive.org/details/diaryofljohnnutt01nutt/page/n135/mode/2up}{Here}
					% \item[Analysis] \hyperref[XX]{Here}
				\end{description}
				
				\begin{description}
					\item[Background]
				\end{description}
				
					\phantom{.}
				
				\begin{description}
					\item[Original Source]
				\end{description}
				
					L. John Nuttall's typescript diary, volume 1, pages 18--21 (starting at PDF 137)
					% brief description of original source material, with reference to JSP or somewhere else for more info
					% Background material: it might be helpful to have a note that says whether a given document was presented as a speech, was written in a journal, etc.
					% Also a comment here about how reliable the source might be, or what shape it is in, if there are large omissions or parts that are hard to understand, and so on.
					% Also a comment about when I transcribed it, whether I've checked it more than once, and so on
				
				\begin{description}
					\item[Text]
				\end{description}
				
					At Temple...Prest Young was filled with the spirit of God \& revelation \& said when we got our washings and anointings under the hands of the Prophet Joseph at Nauvoo we had only one room to work in with the exception of a little side room or office were we were washed and anointed had our garments placed upon us and received our New Name. and after he had performed these ceremonies. he gave the Key Words signs, togkens [sic] and penalties. then after we went into the large room over the store in Nauvoo. Joseph divided up the room the best that he could hung up the veil, marked it gave us our instructions as we passed along from one department to another giving us signs. tokens. penalties with the Key words pertaining to those signs and after we had got through. Bro Joseph turned to me (Prest B Young) and said Bro Brigham this is not arranged right but we have done the best we could under the circumstances in which we are placed, and I...wish you to take this matter in hand and organize and systematize all these ceremonies with the signs. tokens penalties and Key words I did so and each time I got something more so that when we went through the Temple at Nauvoo I understood and Knew how to place them there. we had our ceremonies pretty correct ---
		
					In the creation the Gods entered into an agreement about forming this earth. \& putting Michael or Adam upon it. these things of which I have been speaking are what are termed the mysteries of godliness but they will enable you to understand the expression of Jesus made while in Jerusalem. This is life eternal that they might know thee the only true God and Jesus Christ whom thou hast sent. We were once acquainited [acquainted] with the Gods \& lived with them but we had the privilege of taking upon us flesh that the spirit might have a house to dwell in. we did so and forgot all and came into the world not recollecting anything of which we had previously learned. We have heard a great deal about Adam and Eve. how they were formed \&c some think he was made like an adobie and the Lord breathed into him the breath of life. for we read ``from dust thou art and unto dust shalt thou return'' Well he was made of the dust of the earth but not of this earth. he was made just the same way you and I are made but on another earth. Adam was an immortal being when he came. on this earth he had lived on an earth similar to ours he had received the Priesthood and the Keys thereof. and had been faithful in all things and gained his resurrection and his exaltation and was crowned with glory immortality and eternal lives and was numbered with the Gods for such he became through his faithfulness. and had begotten all the spirit that was to come to this earth. and Eve our common Mother who is the mother of all living bore those spirits in the celestial world. and when this earth was organized by Elohim, Jehovah \& Michael who is Adam our common Father. Adam \& Eve had the privilege to continue the work of Progression. consequently came to this earth and commenced the great work of forming tabernacles for those spirits to dwell in. and when Adam and those that assisted him had completed this Kingdom our earth he came toil. and slept and forgot all and became like an Infant child. it is said by Moses the historian that the Lord caused a deep sleep to come upon Adam and took from his side a rib and formed the woman that Adam called Eve---this should be interpreted that the Man Adam like all other Men had the seed within him to propagate his species. but not the Woman. she conceives the seed but she does not produce it. consequently she was taken from the side or bowels of her father. this explains the mystery of Mose's dark sayings in regard to Adam and Eve. Adam \& Eve when they were placed on this earth were immortal beings with flesh. bones and sinues. but upon partaking of the fruits of the earth while in the garden and cultivating the ground their bodies became changed from immortal to mortal beings with the blood coursing through their veins as the action of life---Adam was not under transgression until affter he partook of the forbidden fruit this was necessary that they might be together that mian might be. the woman was found in transgression not the Man--- Now in the law of Sacrifice we have the promise of a Savior and man had the privilege and showed forth his obedience by offering of the first fruits of the earth and the firstlings of the flocks--- this as a showing that Jesus would come and shed his blood.
					
					Father Adam's oldest son (Jesus the Saviour) who is the heir of the family is Father Adams first begotten in the spirit World. who according to the flesh is the only begotten as it is written. (In his divinity he having gone back into the spirit World. and come in the spirit to Mary and she conceived for when Adam and Eve got through with their Work in this earth. they did not lay their bodies down in the dust, but returned to the spirit World from whence they came.
					
					I felt myself much blessed in being \sout{peraitted} [sic] permitted to Associate with such men and hear such instructions as they savored of life to me---